\documentclass[11pt]{article}
\usepackage[margin=1in]{geometry}
\usepackage{amsmath,amssymb,amsthm}
\usepackage{hyperref}
\hypersetup{hypertexnames=false}
\usepackage{microtype}
\usepackage{booktabs}

\newcommand{\R}{\mathbb{R}}
\newcommand{\eps}{\varepsilon}
\newcommand{\cA}{\mathcal{A}}
\newcommand{\cR}{\mathcal{R}}
\newcommand{\cX}{\mathcal{X}}
\newcommand{\cD}{\mathcal{D}}
\newcommand{\dd}{\mathfrak{d}}
\newcommand{\qq}{\mathfrak{q}}
\newcommand{\sbar}{\bar{s}}
\newcommand{\ybar}{\bar{y}}
\newcommand{\bis}{\operatorname{bis}}
\newcommand{\cl}{\operatorname{cl}}

\newtheorem{theorem}{Theorem}
\newtheorem{lemma}[theorem]{Lemma}
\newtheorem{proposition}[theorem]{Proposition}
\newtheorem{corollary}[theorem]{Corollary}
\theoremstyle{remark}
\newtheorem{remark}[theorem]{Remark}


\title{Distinct distances in general position:\\
       an explicit improvement of the bound $n/3$}
\author{Adam McKenna\thanks{\href{mailto:adam-apple@flounder.net}{adam-apple@flounder.net}}}
\date{October 2026}

\begin{document}
\maketitle

\begin{abstract}
A planar point set is in general position if no three of its points are
collinear and no four are concyclic. For such a set of $n$ points, every
circle centred at a point of the set holds at most three other points, so
the set determines at least $(n-1)/3$ distinct distances. I prove that every
such set with $n \ge 10^4$ determines more than $(1/3 + 1/40)\,n$ distinct
distances. The proof normalizes the set into a thin annulus, derives a
family of linear inequalities for the distribution of pairwise distances,
and closes with an exact rational Farkas certificate over $60$ cases. One
of the inequalities comes from a lemma on the pairs of points that lie
on a common perpendicular bisector. The certificate has been checked by an
independent exact verifier, and the complete proof, the certificate
included, has been checked by the Lean~4 kernel.
\end{abstract}

\section{Introduction}

For a finite set $P \subset \R^2$ let $D(P)$ be the number of distinct
distances $|pq|$, $p \ne q \in P$. We say that $P$ is \emph{in general
position} if no three points of $P$ are collinear and no four points of $P$
are concyclic. Erd\H{o}s asked whether every $n$-point set in general
position determines a number of distinct distances that grows faster than
linearly in $n$ \cite[question~(9)]{E75}; this is Problem~98 on the
Erd\H{o}s problems website \cite{EPweb}, where it is listed as open.

If no four points of $P$ are concyclic, then every circle centred at a point
$p \in P$ contains at most three points of $P$. The $n-1$ other points of
$P$ lie on the circles centred at $p$ whose radii are the distances in $P$,
so $3D(P) \ge n-1$. Brass, Moser and Pach record this bound and state that
no improvement of the form $(1/3+\eps)n$ is known, both for sets with no
four points on a circle and for sets in general position
\cite[\S5.5]{BMP05}. For sets with no three points collinear, Erd\H{o}s
reports an argument of Szemer\'edi that gives a point with more than
$\lfloor n/3\rfloor$ distances from it \cite{E75}. I am not aware of a
bound of the form $(1/3+\eps)n$ for sets in general position. Brass, Moser and Pach
list the upper bound $n e^{O(\sqrt{\log n})}$ for such sets, due to Erd\H{o}s,
F\"uredi, Pach and Ruzsa \cite{EFPR93}. Without the general position
hypothesis the minimum lies between $cn/\log n$, by the theorem of Guth and
Katz \cite{GK}, and $Cn/\sqrt{\log n}$, attained by the square lattice.

The target of this paper is the following statement.

\begin{theorem}\label{thm:main}
Let $P \subset \R^2$ be a set of $n \ge 10^4$ points, no three collinear and
no four concyclic. Then
\[
  D(P) \;>\; \Bigl(\frac13 + \frac1{40}\Bigr)\, n .
\]
\end{theorem}

The proof is by contradiction: it assumes a set of
$n \ge 10^4$ points in general position with $D(P) \le (1/3+1/40)n$ and
derives a contradiction. The constant $1/40$ and the threshold $10^4$ are
explicit, and both come from the certificate of Section~\ref{sec:cert}.

\paragraph{Method.} The proof has four parts.
\begin{enumerate}
\item A counting identity (Lemma~\ref{lem:apex}) relates the deficit
  $Q = 3D - n + 1$ to the number of pairs $\{a,b\}$ whose perpendicular
  bisector contains fewer than two points of $P$.
\item A normalization (Section~\ref{sec:annulus}): after a translation and a
  scaling, if $D \le (1/3+\delta)n$ then every point of $P$ lies in the
  annulus $1 - 12\delta - 4/n \le |p|^2 \le 1$. The width of this annulus
  is controlled by $Q$.
\item Linear inequalities for the distance profile (Sections
  \ref{sec:profile}--\ref{sec:J}): chord bounds for three points on a
  circle inside the annulus, pair counts, a second-moment identity, a cap
  bound, apex bands, and a lemma (Lemma~\ref{lem:J}) on the map that
  sends a pair to the pair of points on its perpendicular bisector.
\item A linear program on a grid of $1283$ radii, split into $60$ cases by
  the mean $\ybar$ of $|p|^2 - \sbar^2$, with an exact rational Farkas certificate for
  each case (Section~\ref{sec:cert}).
\end{enumerate}
Part~4 uses a computer. The certificate is a list of nonnegative rational
multipliers for explicit inequalities; checking it needs only exact
rational arithmetic. Section~\ref{sec:verify} describes the two checks that
were done: an exact verifier that is independent of the program that found
the multipliers, and a Lean~4 formalization of the whole proof in which the
kernel checks the certificate.

\paragraph{Remark on weaker constants.} With fewer inequalities the same
method gives weaker constants without a certificate file: Lemma~\ref{lem:J}
at one radius, Lemma~\ref{lem:P}(b) with a uniform chord bound,
Lemma~\ref{lem:Apex} at $\lambda = \sqrt3$ and Lemma~\ref{lem:M} combine in
$15$ explicit rows to $D(P) > (1/3 + 1/60)n$ for $n \ge 10^4$. That
derivation is not reproduced here.

\section{Notation and the apex identity}\label{sec:setup}

Throughout, $P \subset \R^2$ is a set of $n \ge 2$ points in general
position, $R$ is the set of distances determined by $P$, and $D = |R|$.
For $p \in P$ and $r > 0$ let $C(p,r)$ be the circle with centre $p$ and
radius $r$, and
\[
  d_r(p) = |P \cap C(p,r)|, \qquad \operatorname{def}_r(p) = 3 - d_r(p) .
\]
Four points of $P$ on $C(p,r)$ would be concyclic, so $d_r(p) \le 3$ and
$\operatorname{def}_r(p) \ge 0$. Every $q \in P \setminus \{p\}$ lies on
exactly one circle $C(p,r)$ with $r \in R$, so
$\sum_{r \in R} d_r(p) = n-1$ and
\begin{equation}\label{eq:Q}
  \sum_{r \in R} \operatorname{def}_r(p) \;=\; 3D - n + 1 \;=:\; Q
  \qquad \text{for every } p \in P .
\end{equation}
In particular $Q \ge 0$, which is the bound $D \ge (n-1)/3$.

For distinct $a, b \in P$ let
$\alpha(a,b) = \#\{z \in P\setminus\{a,b\} : |za| = |zb|\}$, the number of
apexes of isosceles triangles with base $ab$. The apexes lie on the
perpendicular bisector $\bis(a,b)$, a line, and no three points of $P$ are
collinear, so $\alpha(a,b) \le 2$. Let $N = \binom n2$ and let $x_k$ be the
number of unordered pairs with $\alpha = k$, so $N = x_0 + x_1 + x_2$.

\begin{lemma}[apex identity]\label{lem:apex}
Let $V = \sum_{p \in P}\sum_{r \in R} v(d_r(p))$ with $v(0) = 3$,
$v(1) = 1$, $v(2) = v(3) = 0$. Then
\[
  \sum_{\{a,b\}} \bigl(2 - \alpha(a,b)\bigr) + V \;=\; nQ ,
\]
the sum running over the $N$ unordered pairs. Consequently
$N - x_2 \le nQ$, and for every set $\Phi$ of ordered pairs of distinct
points that is closed under $(a,b) \mapsto (b,a)$,
\begin{equation}\label{eq:B2}
  \sum_{(a,b) \in \Phi} \bigl(2 - \alpha(a,b)\bigr) \;\le\; 2nQ .
\end{equation}
\end{lemma}

\begin{proof}
Counting the triples ($z$, unordered pair $\{a,b\} \subset C(z,r)$) gives
$\sum_{\{a,b\}} \alpha(a,b) = \sum_{z}\sum_{r\in R}\binom{d_r(z)}{2}$, and
$2N = \sum_z \sum_{r \in R} d_r(z)$. With \eqref{eq:Q},
\[
  nQ - \sum_{\{a,b\}}(2-\alpha) = \sum_z \sum_{r \in R}
  \Bigl[(3 - d) - d + \tbinom d2\Bigr]_{d = d_r(z)},
\]
and $3 - 2d + \binom d2$ equals $3, 1, 0, 0$ for $d = 0, 1, 2, 3$. For the
consequences: $2 - \alpha \ge 0$ and $V \ge 0$, so
$N - x_2 = x_0 + x_1 \le 2x_0 + x_1 = \sum (2 - \alpha) \le nQ$; and the
sum in \eqref{eq:B2} is at most the sum over all ordered pairs, which is
$2\sum_{\{a,b\}}(2-\alpha) \le 2nQ$.
\end{proof}

\paragraph{The map $\kappa$.} Let $\cD$ be the set of unordered pairs $f$
with $\alpha(f) = 2$, so $|\cD| = x_2$. For $f \in \cD$ let $\kappa(f)$ be
the pair of points of $P$ on $\bis(f)$.

\begin{lemma}\label{lem:kappa}
$\kappa$ is injective on $\cD$.
\end{lemma}

\begin{proof}
Let $f \ne g \in \cD$ with $\kappa(f) = \kappa(g) = \{z, z'\}$, and let
$\ell$ be the line $zz'$. Then $\bis(f) = \bis(g) = \ell$, so with $\sigma$
the reflection in $\ell$ we have $f = \{a, \sigma a\}$ and
$g = \{b, \sigma b\}$. If $f$ and $g$ share a point, then $f = g$; so they
are disjoint. If $a$ and $b$ lie on one line perpendicular to $\ell$, the
four points $a, \sigma a, b, \sigma b$ are collinear, which general position
excludes. Otherwise $\bis(a,b)$ is not parallel to $\ell$ and meets $\ell$
at a point $c$, and $|ca| = |cb| = |c\sigma a| = |c\sigma b|$: four points
of $P$ on a circle, which general position also excludes.
\end{proof}

\section{Normalization into an annulus}\label{sec:annulus}

Translate $P$ so that $\sum_{p} p = 0$ and put $S = \sum_p |p|^2$.

\begin{lemma}\label{lem:rho}
$\rho_A^2 := \bigl(3\sum_{r \in R} r^2 - S\bigr)/n$ is positive.
\end{lemma}

\begin{proof}
Since $\sum_q q = 0$, $\sum_q |p-q|^2 = n|p|^2 + S$ for each $p$, so
$\sum_{p,q}|p-q|^2 = 2nS$. Also
$\sum_q |p-q|^2 = \sum_{r\in R} d_r(p) r^2 \le 3\sum_{r\in R} r^2$. Summing
over $p$ gives $2nS \le 3n\sum_r r^2$, so $3\sum_r r^2 - S \ge S$, and
$S > 0$ because $P$ has at least two points and centroid $0$.
\end{proof}

Scale $P$ by $1/\rho_A$. Scaling preserves general position and $D$, and
after it
\begin{equation}\tag{N1}
  3\sum_{r\in R} r^2 = n + S .
\end{equation}
For $p \in P$ put $x_p = n^{-1}\sum_{r \in R}\operatorname{def}_r(p)\,r^2
\ge 0$.

\begin{lemma}\label{lem:annulus}
(N2) $|p|^2 = 1 - x_p$ for every $p \in P$. (N3) $|p| \le 1$, every
distance in $P$ is at most $2$, and $R \subset (0,2]$. (N4) If
$D \le (1/3+\delta)n$, then $x_p \le 12\delta + 4/n$ for every $p$.
\end{lemma}

\begin{proof}
(N2): $n|p|^2 + S = \sum_q|p-q|^2 = \sum_r d_r(p) r^2 =
3\sum_r r^2 - n x_p = n + S - n x_p$ by (N1). (N3) follows from
$x_p \ge 0$. (N4): by (N3) every $r \in R$ satisfies $r^2 \le 4$, so
$x_p \le 4Q/n$ by \eqref{eq:Q}, and $Q = 3D - n + 1 \le 3\delta n + 1$.
\end{proof}

\paragraph{The instance.} From now on assume, for a contradiction, that
$n \ge n_0 = 10^4$ and $D \le (1/3 + \delta)n$ with $\delta = 1/40$. Put
$\bar x = 12\delta + 4/n_0 = 3/10 + 4/10^4$ and
\[
  \sbar = \frac{41821}{50000}, \qquad
  \sbar^2 = \frac{1748996041}{2500000000} \le 1 - \bar x , \qquad
  \omega := 1 - \sbar^2 .
\]
By Lemma~\ref{lem:annulus}, every point of $P$ satisfies
\begin{equation}\label{eq:annulus}
  \sbar^2 \le |p|^2 \le 1 .
\end{equation}
For $p \in P$ put $y_p = |p|^2 - \sbar^2 \in [0,\omega]$ and
$\ybar = n^{-1}\sum_p y_p$.

\section{The distance profile}\label{sec:profile}

For $\rho \ge 0$ let $N_p(\rho) = \#\{q \in P : |pq| \le \rho\}$ ($p$
included) and $R(\rho) = \#\{r \in R : r \le \rho\}$. Put
\[
  \dd = \frac Dn, \quad \eps = \frac1n, \quad \qq = \frac Qn = 3\dd - 1 + \eps,
  \quad \cR(\rho) = \frac{R(\rho)}n, \quad
  \cA(\rho) = \frac1{n^2}\sum_p \bigl(N_p(\rho) - 1\bigr),
\]
and $\cX(\rho) = n^{-2}\sum_p\sum_{r\in R,\, r\le\rho}\operatorname{def}_r(p)$.
So $n^2\cA(\rho)$ is the number of ordered pairs of distinct points at
distance at most $\rho$, and $\#\{\text{unordered pairs } e : |e| \le \rho\}
= n^2\cA(\rho)/2$. Since $\sum_{r\le\rho} d_r(p) = N_p(\rho) - 1$, we have
$\cX = 3\cR - \cA$.

\paragraph{Facts.} (F1) $\cR$, $\cA$ and $\cX$ are nondecreasing,
$0 \le \cX \le \qq$ and $\cR \le \dd$. (F2) $\cR(2) = \dd$ and
$\cA(2) = 1 - \eps$, by (N3). (F3) $1/3 - 1/(3n_0) \le \dd \le 1/3 + \delta$
and $0 < \eps \le 1/n_0$.

\paragraph{Chords.} For $\tau > 0$ and an interval $I$ let $C(I,\tau)$ be
the statement: for every $p \in P$ and every $r \in R\cap I$ with
$d_r(p) = 3$, two of the three points of $P$ on $C(p,r)$ are at distance at
most $\tau$.

\begin{lemma}[same-side chord]\label{lem:C}
Let $p \in P$, $a = |p|$, $r > 0$,
$L = (\sbar^2 - a^2 - r^2)/(2ar)$, $w = (1-\sbar^2)/(2ar)$, let $\cl$ be
the clamp to $[-1,1]$ and $f = \arccos(\cl L) - \arccos(\cl(L+w)) \in
[0,\pi]$. Then any three points of $P$ on $C(p,r)$ contain two at squared
distance at most $2r^2(1 - \cos f)$.
\end{lemma}

\begin{proof}
Rotate so that $p = (a,0)$ and write $q = p + r(\cos\varphi, \sin\varphi)$.
Then $|q|^2 = a^2 + r^2 + 2ar\cos\varphi$, so by \eqref{eq:annulus} a point
$q \in P\cap C(p,r)$ has $\cos\varphi \in [L, L+w]$. Two of three points
lie in the same closed half $\varphi \in [0,\pi]$ or $\varphi \in
[\pi,2\pi]$; reflecting both in the first axis if needed (this preserves
their distance), both lie in $[0,\pi]$. Since $\cos$ is a decreasing
bijection $[0,\pi] \to [-1,1]$, both angles lie in
$[\arccos \cl(L+w), \arccos \cl L]$, so they differ by at most $f$, and the
two points are at distance $2r\sin(|\Delta\varphi|/2) \le 2r\sin(f/2)$.
\end{proof}

\begin{lemma}[box bound]\label{lem:Cbox}
Let $\sbar \le a_0 < a_1 \le 1$ and $0 < r_0 < r_1$. Put
$L_{\mathrm{lo}} = (\sbar^2 - a_1^2 - r_1^2)/(2a_0r_0)$,
$L_{\mathrm{hi}} = (\sbar^2 - a_0^2 - r_0^2)/(2a_1r_1)$,
$w_{\mathrm{hi}} = (1-\sbar^2)/(2a_0r_0)$,
\[
  \Lambda = \{L_{\mathrm{lo}}, L_{\mathrm{hi}}\} \cup
  \bigl(\{-1,\, 1 - w_{\mathrm{hi}}\} \cap [L_{\mathrm{lo}}, L_{\mathrm{hi}}]\bigr),
\]
and let $\beta \le \min_{L \in \Lambda}\bigl[\cl L\cdot \cl(L+w_{\mathrm{hi}})
+ \sqrt{(1 - (\cl L)^2)(1 - (\cl(L+w_{\mathrm{hi}}))^2)}\bigr]$. Then for every
$a \in [a_0,a_1]$ and $r \in [r_0, r_1]$ the bound of Lemma~\ref{lem:C} is at
most $2r_1^2(1-\beta)$. Also, for every $r \in (0, r_1]$, any two points of
$C(p,r)$ are at squared distance at most $4r_1^2$.
\end{lemma}

\begin{proof}
The numerator $\sbar^2 - a^2 - r^2$ is negative and the denominator $2ar$ is
positive, so $L \in [L_{\mathrm{lo}}, L_{\mathrm{hi}}]$ and
$w \le w_{\mathrm{hi}}$. For fixed $L$, $f(L,w)$ is nondecreasing in $w$.
For fixed $w$, $f(\cdot,w)$ is continuous and is: $0$ if $L + w \le -1$;
increasing when $L < -1 < L + w \le 1$; on $-1 \le L$, $L + w \le 1$ its
derivative $-1/\sqrt{1-L^2} + 1/\sqrt{1-(L+w)^2}$ is $\le 0$ for
$L \le -w/2$ and $\ge 0$ for $L \ge -w/2$, so it decreases and then
increases; decreasing when $-1 \le L \le 1 < L+w$; $0$ if $L > 1$; and
$\pi$ on $[1-w, -1]$ when $w > 2$. When $w \le 2$ we have
$-1 \le -w/2 \le 1 - w$, so the only interior local maxima are at $-1$ and
$1-w$. Hence the maximum of $f(\cdot, w_{\mathrm{hi}})$ on
$[L_{\mathrm{lo}}, L_{\mathrm{hi}}]$ is attained at a point of $\Lambda$.
Finally $\cos(\arccos u - \arccos v) = uv + \sqrt{1-u^2}\sqrt{1-v^2}$, so
$\cos f \ge \beta$ on the box, and $2r^2(1-\cos f) \le 2r_1^2(1-\beta)$.
The last claim holds because $C(p,r)$ has diameter $2r$.
\end{proof}

In the certificate, every statement $C(I,\tau)$ that a row needs is proved
by covering $[\sbar, 1] \times \overline I$ with finitely many boxes, on
each of which Lemma~\ref{lem:Cbox} gives a bound at most $\tau^2$, with
square roots rounded so that $\beta$ is a lower bound; a box with $r_0 = 0$
uses the bound $4r_1^2$. When no such cover exists for an interval, the
corresponding row is not used.

\section{Pair counts, moments and a cap}\label{sec:rows}

\begin{lemma}[pair counts]\label{lem:P}
Let $\rho > 0$ and $\rho/2 < r' \le 2$, and assume $C((\rho/2, r'], \rho)$.
Then
\begin{align*}
  &\text{(a)}\quad \cA(\rho) \ge \cA(r') - 2\cR(r'),\\
  &\text{(b)}\quad \cA(\rho) \ge \cA(\rho/2) + \cA(r') - \cR(\rho/2) - 2\cR(r'),\\
  &\text{(c)}\quad \cA(\rho) \ge \cA(r') - \cA(\rho/2) - 2\cR(r') + 2\cR(\rho/2).
\end{align*}
\end{lemma}

\begin{proof}
Let $I_p = \sum_{r \in R}\#\{\{q,q'\} \subset P\cap C(p,r) : |qq'| \le \rho\}$.
An unordered pair $\{q,q'\}$ with $|qq'| \le \rho$ is counted once for each
apex $p$, that is $\alpha(q,q') \le 2$ times, so
$\sum_p I_p \le 2\#\{\{q,q'\}: |qq'| \le \rho\} = n^2\cA(\rho)$. A radius
$r \le \rho/2$ contributes $\binom{d_r(p)}2$ to $I_p$, because all pairs on
$C(p,r)$ are within $2r \le \rho$. A radius $r \in (\rho/2, r']$ contributes
at least $[d_r(p) = 3] \ge d_r(p) - 2$, by the hypothesis. For
$d \in \{0,1,2,3\}$, $\binom d2 \ge d - 2$, $\binom d2 \ge 2d-3$ and
$\binom d2 \ge 0$. Dropping the radii above $r'$ and using
$\sum_{r \le y} d_r(p) = N_p(y) - 1$, the three choices give, for each $p$,
$I_p \ge N_p(r') - 1 - 2R(r')$, $I_p \ge (N_p(\rho/2) - 1) + (N_p(r') - 1) -
R(\rho/2) - 2R(r')$, and $I_p \ge N_p(r') - N_p(\rho/2) - 2R(r') +
2R(\rho/2)$. Sum over $p$ and divide by $n^2$.
\end{proof}

\begin{lemma}[apex range]\label{lem:Apex}
Let $0 < \lambda < 2$ with $\lambda^2 \ge 4\omega$, put
$\mathrm{lo}(\lambda) = \sbar^2 + 1 - \sbar\sqrt{4-\lambda^2}$ and
$\mathrm{hi}(\lambda) = 2 + \sqrt{4 - \lambda^2}$, take
$0 < w^2 \le \mathrm{lo}(\lambda)$ and $\mathrm{hi}(\lambda) \le h^2 \le 4$,
and assume $C((w,2],\lambda)$. Then
\[
  2\bigl(1-\eps-\cA(\lambda)\bigr) - \tfrac43\bigl(\cA(h) - \cA(w)\bigr)
  \;\le\; 2\qq .
\]
\end{lemma}

\begin{proof}
(i) Let $z \notin \{a,b\}$ with $|za| = |zb| = r$ and $|ab| > \lambda$. Put
$u = |a|^2 + |b|^2$ and $\zeta = |z|$. The parallelogram law gives
$|a+b|^2 = 2u - |ab|^2 < 2u - \lambda^2$, and
$2r^2 = |z-a|^2 + |z-b|^2 = 2\zeta^2 + u - 2\langle z, a+b\rangle$. Hence
$\Psi^- < 2r^2 < \Psi^+$ with
$\Psi^\pm = 2\zeta^2 + u \pm 2\zeta\sqrt{2u - \lambda^2}$. By
\eqref{eq:annulus}, $\zeta \in [\sbar,1]$ and $u \in [2\sbar^2, 2]$. Then
$\partial_\zeta \Psi^- \ge 4\sbar - 2\sqrt{4 - \lambda^2} \ge 0$ and, at
$\zeta = \sbar$, $\partial_u\Psi^- = 1 - 2\sbar/\sqrt{2u-\lambda^2} \le 0$
for $u \in (\lambda^2/2, 2]$, both because $4\sbar^2 \ge 4 - \lambda^2$;
the actual $u$ lies in $(\lambda^2/2, 2]$, and $\Psi^-$ is continuous on
$[\lambda^2/2, 2]$. So
$\Psi^- \ge \Psi^-(\sbar, 2) = 2\,\mathrm{lo}(\lambda)$, and
$\Psi^+ \le 4 + 2\sqrt{4-\lambda^2} = 2\,\mathrm{hi}(\lambda)$.

(ii) Let $\Phi$ be the set of ordered pairs $(a,b)$ with $|ab| > \lambda$;
$|\Phi| = n^2(1 - \eps - \cA(\lambda))$. Each apex $z$ of $(a,b) \in \Phi$
lies at a distance $r$ from $a$ with $\mathrm{lo} < r^2 < \mathrm{hi}$ by
(i), so $w < r \le h$. A circle $C(z,r)$ with $d$ points of $P$ carries at
most $\frac43 d$ ordered pairs of $\Phi$: none for $d \le 1$, two for
$d = 2$, and for $d = 3$ at most four, because $C((w,2],\lambda)$ gives a
pair within $\lambda$. Hence $\sum_\Phi \alpha \le \frac43\sum_z\#\{q :
w < |zq| \le h\} = \frac43 n^2(\cA(h) - \cA(w))$. Insert this into
\eqref{eq:B2}, $2|\Phi| - \sum_\Phi\alpha \le 2nQ$, and divide by $n^2$.
\end{proof}

\begin{lemma}[averaged cap]\label{lem:K}
Let $\rho_e^2 = 2\sbar^2 + 2\sbar$. For $0 < g_1 < \dots < g_m \le \rho_e$,
$\sum_{i<m} \cA(g_i)(g_{i+1}^2 - g_i^2) \le 2\sbar$.
\end{lemma}

\begin{proof}
Fix $z \in P$, $\zeta = |z| \in [\sbar, 1]$, and put
$F(q) = 1 + \langle q, z/\zeta\rangle$. By (N3), $F \ge 0$ on $P$, and
$\sum_q F(q) = n$ because the centroid is $0$. Since
$\langle q, z/\zeta\rangle = (|q|^2 + \zeta^2 - |q-z|^2)/(2\zeta)$,
\eqref{eq:annulus} gives $F(q) \ge \varphi_\zeta(|q-z|)$ with
$\varphi_\zeta(t) = 1 + (\sbar^2 + \zeta^2 - t^2)/(2\zeta)$. Its
derivative in $\zeta$ is $\frac12 - (\sbar^2 - t^2)/(2\zeta^2) \ge 0$
for $\zeta \ge \sbar$, so $\varphi_\zeta \ge \varphi_{\sbar} =
(\rho_e^2 - t^2)/(2\sbar)$. Therefore
$n \ge \sum_q \max(0, \varphi_{\sbar}(|qz|)) =
\frac1{\sbar}\int_0^{\rho_e} t\,N_z(t)\,dt$. Summing over $z$,
$\int_0^{\rho_e} t\,(n^2\cA(t) + n)\,dt \le \sbar n^2$, so
$\int_0^{\rho_e} t\cA(t)\,dt \le \sbar$. Since $\cA$ is nondecreasing, the
left side of the claim is at most twice this integral.
\end{proof}

\begin{lemma}[second moment]\label{lem:M}
$\int_0^2 \rho\,\bigl(\cA(\rho) + 2\cX(\rho)\bigr)\,d\rho = 1 + 4\qq - 2\eps$.
The integrand is nondecreasing by (F1), so on a partition
$0 = g_0 < \dots < g_m = 2$ the sum
$\sum_i \frac12(g_{i+1}^2 - g_i^2)(\cA + 2\cX)(g_i)$ is at most
$1 + 4\qq - 2\eps$, and the same sum with $(\cA + 2\cX)(g_{i+1})$ is at
least $1 + 4\qq - 2\eps$.
\end{lemma}

\begin{proof}
For each ordered pair $(p,q)$, $p = q$ allowed,
$2 - |pq|^2/2 = \int_{|pq|}^2 \rho\,d\rho$ by (N3). Summing and dividing by
$n^2$: the left side becomes $\int_0^2 \rho(\cA + \eps)\,d\rho$, and the
right side becomes $2 - S/n = 1 + \bar x_{\mathrm{avg}}$ with
$\bar x_{\mathrm{avg}} = n^{-1}\sum_p x_p$, by $\sum_{p,q}|p-q|^2 = 2nS$ and
(N2). Finally $\bar x_{\mathrm{avg}} = \int_0^2 \rho^2\,d\cX(\rho) =
4\cX(2) - 2\int_0^2\rho\cX(\rho)\,d\rho$ with $\cX(2) = \qq$ by (N3) and
\eqref{eq:Q}.
\end{proof}

\section{Apex bands}\label{sec:AB}

Fix a rational $\lambda < 2$ such that $C((0,2],\lambda)$ holds. Let $\Phi$
be the set of ordered pairs with $|ab| > \lambda$ and
$T = \{(z,a,b) : (a,b) \in \Phi,\ z \notin\{a,b\},\ |za| = |zb|\}$. Fix
$e_0 < e_1 < \dots < e_m = 2$ with $e_0 \le \lambda/2$, let $T_j$ be the
triples with $e_j < |za| \le e_{j+1}$, and put $W_j = |T_j|/n^2$,
$U_j = n^{-2}\sum_{T_j}(y_a + y_b)$, $V_j = n^{-2}\sum_{T_j} y_z$, and for
$t > 0$, $K^\pm(t) = (4 \pm t)\sbar^2 \pm (4\sbar^2 - \lambda^2)/t$.

\begin{lemma}[apex bands]\label{lem:AB}
\begin{itemize}
\item[(AB0)] $\ybar = \omega - 4\qq + 2\int_0^2 \rho\,\cX(\rho)\,d\rho$.
\item[(AB1)] $\sum_j W_j \ge 2(1 - \eps - \cA(\lambda)) - 2\qq$.
\item[(AB2)] $W_j \le \frac43\bigl(\cA(e_{j+1}) - \cA(e_j)\bigr)$.
\item[(AB3)] $(1 + 2/t)U_j + (2+t)V_j \ge (2e_j^2 - K^+(t))W_j$.
\item[(AB4)] $(2/t - 1)^+U_j + (t-2)^+V_j \ge (K^-(t) - 2e_{j+1}^2)W_j$.
\item[(AB5)] $\sum_j U_j \le 12\,\ybar\,(\dd - \cR(\lambda))$ and, for
  every $h \le \lambda/2$, $\sum_j V_j \le 4\,\ybar\,(\dd - \cR(h))$.
\end{itemize}
\end{lemma}

\begin{proof}
Every triple of $T$ has $2|za| \ge |ab| > \lambda$ and $|za| \le 2$, so it
lies in exactly one band. (AB0): $n^{-1}\sum_p|p|^2 = 1 - \bar
x_{\mathrm{avg}}$ by (N2), and $\bar x_{\mathrm{avg}} = 4\qq -
2\int_0^2\rho\cX$ as in the proof of Lemma~\ref{lem:M}. (AB1):
$|T| = \sum_\Phi\alpha \ge 2|\Phi| - 2nQ$ by \eqref{eq:B2}. (AB2): fix $z$
and a radius $r$ in the band; $C(z,r)$ carries at most $\frac43 d_r(z)$
ordered pairs of $\Phi$, as in Lemma~\ref{lem:Apex}(ii), now using
$C((0,2],\lambda)$; sum over the radii in the band and over $z$.

For (AB3) and (AB4) take $(z,a,b) \in T$, $r = |za|$, $\ell = |ab|$, and
write $|z|^2 = \sbar^2 + y_z$, $u = |a|^2 + |b|^2 = 2\sbar^2 + Y$ with
$Y = y_a + y_b$, and $\Theta = 2u - \lambda^2 = 4\sbar^2 - \lambda^2 + 2Y$.
As in Lemma~\ref{lem:Apex}, $2r^2 = 2|z|^2 + u - 2\langle z, a+b\rangle$ and
$|a+b|^2 = 2u - \ell^2 < \Theta$, so $\Theta > 0$ and, by
the Cauchy--Schwarz and AM--GM inequalities,
$|2\langle z,a+b\rangle| < 2\sqrt{|z|^2 \Theta} \le t|z|^2 + \Theta/t$.
Hence
\[
  K^-(t) + (1 - 2/t)Y + (2-t)y_z \;<\; 2r^2 \;<\; K^+(t) + (1+2/t)Y + (2+t)y_z .
\]
On $T_j$, $r > e_j$; summing the right inequality over $T_j$ gives (AB3).
On $T_j$, $r \le e_{j+1}$, so $(2/t - 1)Y + (t-2)y_z > K^-(t) - 2e_{j+1}^2$;
replace the coefficients by their positive parts ($Y, y_z \ge 0$) and sum,
which gives (AB4).

(AB5): $\Phi$ is symmetric, so $\sum_T (y_a + y_b) = 2\sum_{(a,b)\in\Phi}
\alpha(a,b)\,y_a \le 4\sum_a y_a M_a$ with $M_a = \#\{b : |ab| > \lambda\}
= \sum_{r > \lambda} d_r(a) \le 3(D - R(\lambda))$. The ordered pairs of
$\Phi$ with apex $z$ lie on circles $C(z,r)$ with $r > \lambda/2 \ge h$,
and by the count in (AB2) there are at most $\frac43\sum_{r>h} d_r(z) \le
4(D - R(h))$ of them. Multiply by $y_z$, sum, and divide by $n^2$.
\end{proof}

\section{The bisector length gap}\label{sec:J}

The inequalities so far use the apex counts $\alpha$ only. The next lemma
uses the image $\kappa(\cD)$: on the annulus, the pair of apexes of a long
pair is either short or long.

\begin{lemma}[bisector length gap]\label{lem:J}
Let $X \subset \R^2$ be finite, no three points collinear and no four
concyclic, with $\sbar^2 \le |p|^2 \le 1$ for all $p \in X$, where
$0 < \sbar < 1$, and let $\omega = 1 - \sbar^2$; define $\cD$, $\kappa$,
$N$ and $x_2$ for $X$ as in Section~\ref{sec:setup}. Fix $\lambda > 0$ with
$h_0 := \omega/(2\lambda) < \sbar$ and put
\[
  \sigma_0 = \sqrt{1 - h_0^2} - \sqrt{\sbar^2 - h_0^2}, \qquad
  \Lambda_0 = 2\sqrt{\sbar^2 - h_0^2} .
\]
\begin{itemize}
\item[(i)] If $f \in \cD$ and $|f| \ge \lambda$, then $|\kappa(f)| \le \sigma_0$
  or $|\kappa(f)| \ge \Lambda_0$.
\item[(ii)] $\#\{e : \lambda \le |e| < \Lambda_0\} \le \#\{e : |e| \le \sigma_0\}
  + (N - x_2)$, where $e$ runs over unordered pairs.
\end{itemize}
\end{lemma}

(No centroid condition is used; the lemma is applied to $X = P$.)

\begin{proof}
(i) For $f = \{a,b\}$, $\bis(f) = \{x : 2\langle x, b - a\rangle = |b|^2 -
|a|^2\}$ lies at distance $h = \bigl||b|^2 - |a|^2\bigr|/(2|ab|) \le
\omega/(2\lambda) = h_0 < \sbar$ from $0$. Write $\bis(f) = \{c + tu : t \in
\R\}$ with $c$ the foot of the perpendicular from $0$ and $|u| = 1$, so
$|c + tu|^2 = h^2 + t^2$. The two points $z = c + tu$, $z' = c + t'u$ of
$\kappa(f)$ satisfy $\sbar^2 \le h^2 + t^2, h^2 + t'^2 \le 1$, so
$\sqrt{\sbar^2 - h^2} \le |t|, |t'| \le \sqrt{1-h^2}$; in particular $t$ and
$t'$ are nonzero. If $tt' > 0$, then $|zz'| = \bigl||t| - |t'|\bigr| \le
\sigma(h) := \sqrt{1-h^2} - \sqrt{\sbar^2 - h^2}$. If $tt' < 0$, then
$|zz'| = |t| + |t'| \ge \Lambda(h) := 2\sqrt{\sbar^2 - h^2}$. On $[0,\sbar)$,
$\sigma'(h) = h\bigl(1/\sqrt{\sbar^2 - h^2} - 1/\sqrt{1 - h^2}\bigr) \ge 0$
and $\Lambda$ is decreasing, so $\sigma(h) \le \sigma_0$ and
$\Lambda(h) \ge \Lambda_0$.

(ii) Let $F = \{f \in \cD : |f| \ge \lambda\}$. By (i) and
Lemma~\ref{lem:kappa}, $|F| \le \#\{|e| \le \sigma_0\} + \#\{|e| \ge
\Lambda_0\}$. Every pair with $|e| \ge \lambda$ that is not in $F$ is
outside $\cD$, so $|F| \ge \#\{|e| \ge \lambda\} - (N - x_2)$. Subtracting
gives (ii). (If $\Lambda_0 \le \lambda$, the left side of (ii) is $0$.)
\end{proof}

\begin{corollary}[row J]\label{cor:J}
In the setting of Section~\ref{sec:annulus}, let $\lambda$ be as in
Lemma~\ref{lem:J}, and let $g_2 < \Lambda_0$ and $g_\sigma \ge \sigma_0$ with
$g_2 > \lambda$. Then $\cA(g_2) - \cA(\lambda) - \cA(g_\sigma) \le 2\qq$.
\end{corollary}

\begin{proof}
$\{\lambda < |e| \le g_2\} \subseteq \{\lambda \le |e| < \Lambda_0\}$ and
$\{|e| \le \sigma_0\} \subseteq \{|e| \le g_\sigma\}$. Apply
Lemma~\ref{lem:J}(ii) and $N - x_2 \le nQ$ (Lemma~\ref{lem:apex}), multiply
by $2/n^2$, and use $\#\{|e|\le\rho\} = n^2\cA(\rho)/2$.
\end{proof}

\begin{remark}
Corollary~\ref{cor:J} adds information to the other families. A closely
related system without it (the same families with $34$ thresholds and a
slightly different set of (AB5)-type rows, built at $\delta = 1/44$ and at
$\delta = 177/7832$) has exact rational feasible points, and these points
violate Corollary~\ref{cor:J} by about $0.2827$ and $0.3045$ (exact
rational evaluation). At $\delta = 1/40$ the system of
Section~\ref{sec:cert} without Corollary~\ref{cor:J} is feasible in
floating-point arithmetic.
\end{remark}

\section{The certificate}\label{sec:cert}

\paragraph{Grid.} Let $G$ be the set of rationals $b/2^j \ge 1/125$
($j \ge 0$), where $b$ runs over the $160$ numbers
$2\cdot 2^{-k/160}$ ($0 \le k < 160$) rounded to six decimals, the number
$\lceil\sqrt3\cdot 10^6\rceil/10^6$, and $2$. Then $|G| = 1283$.

\paragraph{Variables and box.} The unknowns are $\dd$, $\eps$, $\cR(g)$ and
$\cA(g)$ for $g \in G$, and the band quantities $W_j, U_j, V_j$ of
Section~\ref{sec:AB} for $32$ thresholds $\lambda$ (below). The true values
satisfy the box: (F3), $\cR(g) \in [0, 1/3 + \delta]$, $\cA(g) \in [0,1]$,
$W_j \in [0,2]$ and $U_j, V_j \in [0,4]$ (since $|T| = \sum_\Phi\alpha \le
2n^2$, $U_j \le 2\omega W_j$ and $V_j \le \omega W_j$).

\paragraph{Rows.} All rows are linear inequalities in these unknowns whose
coefficients are explicit rationals; each instance carries its parameters,
and every side condition is checked in exact arithmetic.
\begin{enumerate}
\item (F1), (F2) on $G$: monotonicity of $\cR$, $\cA$ and $\cX = 3\cR -
  \cA$ between consecutive grid points, $\cX \ge 0$, $\cX \le \qq$,
  $\cR \le \dd$, $\cR(2) = \dd$, $\cA(2) = 1 - \eps$.
\item Lemma~\ref{lem:P} (a)--(c) for each $\rho \in G$ for which a grid
  point $r' > \rho/2$ with $C((\rho/2, r'], \rho)$ is certified by
  Lemma~\ref{lem:Cbox} ($579$ values of $\rho$; $r'$ is read from the
  certifier's chord table, and $r' = 2$ when the whole range is
  certified); (b) and (c) only when $\rho/2 \in G$.
\item Lemma~\ref{lem:Apex} for every $\lambda \in G$ with
  $4\omega \le \lambda^2 < 4$, $w$ the largest grid point with $w^2 \le
  \mathrm{lo}(\lambda)$ (rounded down) and $h$ the smallest grid point with
  $h^2 \ge \mathrm{hi}(\lambda)$ (rounded up), when $C((w,2],\lambda)$ is
  certified.
\item Lemma~\ref{lem:K} on the grid points up to $\rho_e$ (rounded down),
  and the lower and upper Riemann sums of Lemma~\ref{lem:M} on $G$.
\end{enumerate}
These are $9577$ rows. Then:
\begin{enumerate}\setcounter{enumi}{4}
\item Lemma~\ref{lem:AB}, (AB1)--(AB4), for $32$ thresholds
  $\lambda \in G \cap [1, 2\sbar + 1/20)$, each with $C((0,2],\lambda)$
  certified, chosen greedily upward from $1$ with consecutive thresholds at
  least $1/50$ apart; band edges $e_0 =$ the largest grid point
  $\le \lambda/2$ and $e_k =$ the largest grid point
  $\le \min(e_0 + k/50, 2)$, ending at $2$; and
  $t \in \{(1/50)(5/4)^k : 0 \le k < 25\}$, a row of (AB3) or (AB4) being
  used only when its coefficient of $W_j$ is positive. These are $15182$
  rows. (The certificate file also lists $64$ rows of a weaker form of
  (AB5) with weight zero.)
\item Corollary~\ref{cor:J} for every $\lambda \in G$ with $h_0 < \sbar$:
  $g_2$ is the largest grid point below a rational lower bound of
  $\Lambda_0$, $g_\sigma$ the smallest grid point at or above a rational
  upper bound of $\sigma_0$, and $g_2 > \lambda$. These are $515$ rows.
\end{enumerate}

\paragraph{Case split.} $\ybar \in [0,\omega]$. Split $[0,\omega]$ into the
$60$ intervals $[Y_k, Y_{k+1}] = \omega\,[k/60, (k+1)/60]$. In case $k$
($\ybar \in [Y_k, Y_{k+1}]$) add, for each threshold $\lambda$ with
$h = $ the largest grid point $\le \lambda/2$,
\[
  \textstyle\sum_j U_j \le 12\,Y_{k+1}(\dd - \cR(\lambda)), \qquad
  \sum_j V_j \le 4\,Y_{k+1}(\dd - \cR(h)),
\]
which follow from (AB5) and $\dd - \cR \ge 0$, and the two rows
\[
  Y_k \le \omega - 4\qq + 2\,\mathrm{Up}, \qquad
  \omega - 4\qq + 2\,\mathrm{Lo} \le Y_{k+1},
\]
where $\mathrm{Up}$ and $\mathrm{Lo}$ are the upper and lower Riemann sums
of $\int_0^2\rho\cX$ on $G$ ($\mathrm{Up}$ includes the cell $[0,\min G]$
with the value $\cX(\min G)$, $\mathrm{Lo}$ omits it); they follow from (AB0) because $\cX$ is
nondecreasing and nonnegative. The variable $\ybar$ itself does not occur.

\paragraph{The Farkas step.} Write the rows of case $k$ as $Av \le b$, where
$v$ is the vector of true values. For each $k$ the certificate gives
nonnegative rational multipliers $y$ with
\[
  \min_{v \in \text{box}}\; (y^{\mathsf T}A)\,v - y^{\mathsf T}b \;>\; 0 .
\]
The minimum of a linear function over a box is computed coordinatewise, in
exact arithmetic. But the true $v$ lies in the box and satisfies $Av \le b$,
so $(y^{\mathsf T}A)v \le y^{\mathsf T}b$: a contradiction. Each case uses
$68$ to $130$ instances of Corollary~\ref{cor:J}, and the $60$ cases use
$125126$ multiplier rows in all. The smallest gap is in case $30$
($\ybar \in [0.1502, 0.1552]$):
\[
  \frac{42734354885632410509310622714769366500859}
       {13107200000000000000000000000000000000000000000}
  \;\approx\; 3.2604\cdot 10^{-6}.
\]
Since every case is contradictory and $\ybar$ lies in one of them, the
assumption $D \le (1/3+1/40)n$ with $n \ge 10^4$ is false. This proves
Theorem~\ref{thm:main}. \qed

\section{Verification}\label{sec:verify}

\paragraph{Exact verifier.} The multipliers were found by a floating-point
linear-programming solver and then rounded to rationals; the solver plays no
part in the proof. A separate verifier, written with the Python standard
library only and not importing the program that found the multipliers,
rebuilds every row from its tag, re-checks every side condition (grid
membership of each parameter, the chord covers of Lemma~\ref{lem:Cbox}, the
square-root bounds in Lemma~\ref{lem:Apex} and Corollary~\ref{cor:J}, the
cut of Lemma~\ref{lem:K}, the band edges and the case bounds), and
evaluates the box minimum in exact rational arithmetic for each of the $60$
cases. In two independent audits the check rejected mutated row tags and a
case with the rows of Corollary~\ref{cor:J} removed.

\paragraph{Lean.} The whole proof is formalized in Lean~4 with Mathlib. Each
row family is a Lean lemma about an arbitrary configuration in general
position under rational side conditions; a computable row builder maps a
tag to a rational row, checks the side conditions, and is proved sound. For
each of the $60$ cases a single theorem checks the weighted combination by
kernel reduction (\texttt{decide +kernel}), with the multipliers rounded
to integers at scale $2^{200}$. The final statement,
\texttt{Esgk.numDistances\_gt\_one\_third\_add\_kappa\_image\_gap}, says
that for an injective configuration of $n \ge 10000$ points in general
position, $(1/3 + 1/40)\,n < D$. Its axiom closure consists of the three standard
axioms \texttt{propext}, \texttt{Classical.choice} and \texttt{Quot.sound};
the proof uses no \texttt{sorry} and no \texttt{native\_decide}.

\paragraph{Availability.} The certificate, the exact verifier and the Lean
sources are at
\url{https://github.com/mysticflounder/distinct-distances-general-position}.

\section{Remarks}

\paragraph{Limits of the method.} With the same row families, the
floating-point linear program has feasible points at $\delta = 1/39$ for
$\ybar \in \{0.128, 0.154, 0.180, 0.205\}$. This indicates (floating point
only; EMPIRICALLY VERIFIED in that sense) that these rows alone do not give
$\delta = 1/39$ by this kind of certificate. At $\ybar = 0.154$, three
further families of inequalities (Bochner positivity of the distance
measure, a lower bound for short pairs from the area of the annulus, and a
three-point covering inequality) each cut the feasible point found but
leave the linear program feasible. HEURISTIC: a larger constant needs a
further inequality of the kind of Lemma~\ref{lem:J}.

\paragraph{Pinned distances.} Theorem~\ref{thm:main} bounds the total number
of distances. It gives no bound for the number $t(p)$ of distances from a
single point, and the questions of Erd\H{o}s and Pach whether
$\max_p t(p) \ge (1/3+\eps)n$ or $\sum_p t(p) \ge (1/3+\eps)n^2$ for sets
in general position \cite[Problem~2]{EP90}; I am not aware of an answer
to either. The second question is tied to the number $x_2$ of pairs with
two apexes. Let $c_k$ be the number of pairs $(p,r)$, $r > 0$, with
$d_r(p) = k$. As in the proof of Lemma~\ref{lem:apex},
$2x_0 + x_1 = \sum_{\{a,b\}}(2 - \alpha) =
\sum_{p,r}(d - \binom d2)_{d = d_r(p)} = c_1 + c_2$, and
$\sum_p (3t(p) - (n-1)) = \sum_{p,r:\, d_r(p) \ge 1}(3 - d_r(p)) =
2c_1 + c_2$. Since $x_0 + x_1 \le 2x_0 + x_1 \le 2(x_0 + x_1)$ and
$c_1 + c_2 \le 2c_1 + c_2 \le 2(c_1 + c_2)$, the quantities $N - x_2$ and
$\sum_p (3t(p) - (n-1))$ are within a factor $4$ of each other. So a bound
$x_2 \le (1-c)N$ with $c > 0$ gives a positive answer to the second
question, and a positive answer gives such a bound. Erd\H{o}s and Purdy
\cite{EPu76} remark that, for sets in general position, they cannot even
show that the number of isosceles triangles is $o(n^2)$. The proof above
uses the common set $R$ of distances through \eqref{eq:Q} and does not
bound $x_2$.

\appendix
\section{Status of the ingredients}

\begin{center}
\small
\begin{tabular}{lp{5.2cm}p{3.6cm}l}
\toprule
Obligation & Statement & Label & Evidence\\
\midrule
O1 & Lemma~\ref{lem:apex} (apex identity) & PROVEN (paper, Lean) & \S\ref{sec:setup}\\
O2 & Lemma~\ref{lem:kappa} ($\kappa$ injective) & PROVEN (paper, Lean) & \S\ref{sec:setup}\\
O3 & Lemmas~\ref{lem:rho}, \ref{lem:annulus} (annulus) & PROVEN (paper, Lean) & \S\ref{sec:annulus}\\
O4 & Lemmas~\ref{lem:C}, \ref{lem:Cbox} (chords) & PROVEN (paper, Lean) & \S\ref{sec:profile}\\
O5 & Lemmas~\ref{lem:P}--\ref{lem:M} & PROVEN (paper, Lean) & \S\ref{sec:rows}\\
O6 & Lemma~\ref{lem:AB} (apex bands) & PROVEN (paper, Lean) & \S\ref{sec:AB}\\
O7 & Lemma~\ref{lem:J}, Corollary~\ref{cor:J} & PROVEN (paper, Lean) & \S\ref{sec:J}\\
O8 & the 60 case certificates & PROVEN (exact check, Lean kernel) & \S\ref{sec:cert}, \S\ref{sec:verify}\\
O9 & Theorem~\ref{thm:main} & PROVEN (paper, Lean) & \S\ref{sec:verify}\\
O10 & no certificate at $\delta = 1/39$ from these rows & EMPIRICALLY VERIFIED & floating-point LP only\\
\bottomrule
\end{tabular}
\end{center}

\begin{thebibliography}{9}
\bibitem{BMP05} P.~Brass, W.~Moser, J.~Pach, \emph{Research Problems in
  Discrete Geometry}, Springer, New York, 2005.
\bibitem{E75} P.~Erd\H{o}s, On some problems of elementary and combinatorial
  geometry, \emph{Ann.\ Mat.\ Pura Appl.} (4) \textbf{103} (1975).
\bibitem{EPu76} P.~Erd\H{o}s, G.~Purdy, Some extremal problems in geometry
  IV, \emph{Congr.\ Numer.} \textbf{17} (1976).
\bibitem{EFPR93} P.~Erd\H{o}s, Z.~F\"uredi, J.~Pach, I.~Z.~Ruzsa, The grid
  revisited, \emph{Discrete Math.} \textbf{111} (1993);
  doi:10.1016/0012-365X(93)90155-M.
\bibitem{EP90} P.~Erd\H{o}s, J.~Pach, Variations on the theme of repeated
  distances, \emph{Combinatorica} \textbf{10} (1990).
\bibitem{EPweb} T.~F.~Bloom, Erd\H{o}s Problem \#98,
  \url{https://www.erdosproblems.com/98}, accessed 24 October 2026.
\bibitem{GK} L.~Guth, N.~H.~Katz, On the Erd\H{o}s distinct distances
  problem in the plane, \emph{Ann.\ of Math.} (2) \textbf{181} (2015),
  155--190; arXiv:1011.4105.
\end{thebibliography}

\end{document}
